\section{Erd\H{o}s Problem \#306: independent continuation}
\noindent Date: 23 September 2026. Source versions read in full: \texttt{306.tex} in \texttt{gpt\_pro\_5.4}.

\subsection*{1) OBJECTIVE AND SOURCE AUDIT}
The question concerns every positive rational with squarefree reduced
denominator, using distinct products of two distinct primes. The supplied
graph sum identity is correct but selects the numerator rather than
representing an arbitrary prescribed one. I use the same graph to obtain
necessary cancellation constraints on any putative representation.

\subsection*{2) WORK: WHERE LEAVES ARE ALLOWED}
Suppose
\[
 a/b=\sum_{\{p,q\}\in E(G)}\frac1{pq},\qquad (a,b)=1,
\]
where $G$ is a finite simple graph on its nonisolated prime vertices.
Put $P=\prod_{p\in V(G)}p$. Necessarily $b\mid P$. Multiplying the equation
by $P$ and reducing modulo a vertex prime $p\nmid b$ gives
\[
 \sum_{q\in N_G(p)}q^{-1}\equiv0\pmod p. \tag{1}
\]
Indeed, the left side $aP/b$ is divisible by $p$, as are all edge
contributions not incident with $p$. Dividing the remaining congruence
by $P/p$, which is invertible modulo $p$, proves (1).

In particular every leaf prime must divide $b$: a single inverse cannot
vanish. Hence a representation of an integer has minimum graph degree
at least two. A representation with prime reduced denominator also
cannot have a forest as its support graph. Every nonempty finite forest
has at least two leaves, whereas at most one vertex prime can divide
such a denominator. Thus these representations must use at least one
cycle, not merely a tree of independently introduced primes.

More generally every nonempty tree component requires at least two
distinct prime factors of $b$ among its leaves. If the graph has $t$
tree components, then $2t\le\omega(b)$. This is useful before any
large-integer factorization in a proposed construction.

\subsection*{3) ADVERSARIAL VERIFICATION}
The constraints are necessary only. A cycle does not automatically
represent the required rational. The source's triangle identity passes
this test, but provides no decomposition for arbitrary $a/b$.
Also, repeating the same identity to represent multiples is invalid if
it repeats denominators; simplicity of $G$ records that restriction.

\subsection*{7) FINAL AND REFERENCE}
\textbf{UNRESOLVED.} This attempt supplies structural obstructions, not a
general construction. The indexed official discussion reports a claimed
formalized solution; I have not checked its axiom boundary or proof and
do not adopt that claim. The saved repository status is still open.
Target checked at \url{https://www.erdosproblems.com/306}.

\subsection*{8) COMPLETION ESTIMATE}
\textbf{COMPLETION: 8\%}
This is a subjective estimate of progress toward the entire stated problem, not a probability that the conjecture or an unverified proof is correct.
